% s-133-p133-m002; original PDF 133, 134
\Aphorism{s-133-p133-m002}{188}
In unserem modernsten Empfinden wird die \Sperr{Ehe} wieder zum Ideal, und vielleicht zu einem höheren, mystischer angebeteten als sie es je war. Wir haben den erotischen Liberalismus so satt wie den religiösen, Aphrodite Pandemos so satt wie die „freie Forschung“. Ehedem war man Freigeist über Gott und menschliches Elend; es gab schöne grüne Zeiten, wo man sich für ausbündig tapfer hielt, wenn man der Kirche Unfläthigkeiten sagte. Als diese Themen ausgebeutet waren, machte man \Foreign{esprit} über Weib, Liebe und Ehe, und wieder waren es die lautesten Verunglimpfer, die am wenigsten das Recht hatten mitzureden. Diejenigen, die am Theismus als an einem wirklichen Entwicklungshemmniss gelitten haben, sind ebenso schweigsame Atheisten (oder vielleicht Antitheisten), wie Diejenigen, die am Weibe gescheitert sind, sich zwar vom Weibe, aber ebenfalls von den wüthenden Verleumdungen des „Ewig-Weiblichen“ fernhalten. Die Unbetheiligten, die nur aus zweiter Hand wissen, haben immer Geist, den Geist hinzuzulügen, Pointen zu spitzen, Bosheiten einzufädeln; der wirklich Betroffene, wenn er überhaupt wieder zu sich kommt, steht dann seinem Erlebniss fern genug, um schlicht und ohne Groll darüber zu sprechen. — Nun, wie gesagt, der Geschmack an den \Foreign{rancuniers} und ihren Arten der Polemik ist im Abnehmen; Gott und das Weib werden nicht nur rehabilitirt, sondern auf neue Altäre gestellt. Immer sind es noch Wenige, die ernsthaft, will sagen gleich fern von Compromiss wie von Feindseligkeit, den Versuch machen, es ohne Religion und Liebe auszuhalten.
\par

